\section*{5 A deterministic reduction that works on average}
We will now use the tools developed in Section 3 and Section 4 to give a search to decision reduction that is efficient on functions with few near-optimal formulas.
\setcounter{nttheorem}{32}
\begin{nttheorem}[Author Theorem 33]
There is a deterministic algorithm solving Search-MFSP on inputs of length $N$ given access to an oracle that solves MFSP on instances of length $2N$ that runs in time
\[O(|\Canon{n+1}(f)|^2\cdot N^6\log^2N)\]
where $n=\log N$.
\end{nttheorem}
Before we prove Theorem 33, we state a corollary that follows from the bound on the size of $\Canon{k}(f)$ for a random function given in Proposition 14.
\begin{ntcorollary}[Author Corollary 34]
There is a deterministic algorithm solving Search-MFSP on inputs of length $N$ given access to an oracle that solves MFSP on instances of length $2N$ that runs in time $2^{O(N/\log\log N)}$ on all but a $o(1)$-fraction of instances.
\end{ntcorollary}
\begin{proof}[Author proof of Corollary 34]
This corollary follows by combining the algorithm in Theorem 33 with the recursive algorithm for Search-MFSP given in Theorem 20 and using Proposition 14 to bound $|\Canon{n+1}(f)|$ by $2^{O(N/\log\log N)}$ for a random function.
There is actually a subtlety in appealing to Theorem 20 in that the running time of the algorithm in Theorem 33 has a dependence on the number near-optimal formulas of its input. Hence, we need to argue that when the algorithm in Theorem 20 makes recursive calls to the algorithm in Theorem 20 on functions $g$ other than the original input $f$ that $g$ also has few near-optimal formulas. However, this is not a problem since it is easy to see that if a function $g$ is computed by some gate in an optimal formula for $f$ (as it must be if a recursive call is made to $g$), then the number of near-optimal formulas for $g$ is at most the number of near-optimal formulas for $f$ (since one can create a near optimal formula for $f$ by taking the optimal formula for $f$ that computes $g$ at some gate and replacing the subformula at that gate with a near-optimal formula for $g$).
\end{proof}
\begin{proof}[Author proof of Theorem 33]
We provide the pseudocode of our DecompProblem algorithm in Algorithm 2, which we recommend the reader look at before proceeding.
\begin{quote}\textbf{Algorithm 2. A deterministic search to decision reduction for MFSP whose run time depends on the number of ``near-optimal formulas''.}\par
\textbf{procedure OptimalSubcomputation($f$)}\par
Given the length-$N$ truth table of a function $f$ that takes $n$-inputs with $L(f)\geq 2$, this procedure returns an element $\{g,O,h\}$ of $\Opt(f)$.
\paragraph{Part 1: Building a Candidates list}
Let $\operatorname{allUnknown}:\{0,1\}^n\to\{0,1,?\}$ be given by $\operatorname{allUnknown}(x)=?$ for all $x$. Set $\PC{0}=\{\operatorname{allUnknown}\}$. Set $i=0$.
\textbf{while} $i<N$ \textbf{do}
\begin{enumerate}
\item Set $\PC{i+1}=\varnothing$.
\item For all $\gamma\in\PC{i}$ and for all $b\in\{0,1\}$ do
\begin{enumerate}
\item Let $x^?$ (original line 10) be the lexicographically first input satisfying $\gamma(x^?)=?$.
\item Let $\gamma_b:\{0,1\}^n\to\{0,1,?\}$ be given by $\gamma_b(x)=\begin{cases}b,&\text{if }x=x^?,\\\gamma(x),&\text{otherwise}.\end{cases}$
\item Let $g_{\gamma_b}$ be the (total) function given by $g_{\gamma_b}(x)=\begin{cases}1,&\text{if }\gamma_b(x)=?,\\\gamma_b(x),&\text{otherwise}.\end{cases}$
\item If $L(\Select[f,g_{\gamma_b}])\leq L(f)+n+2$ then add $\gamma_b$ to $\PC{i+1}$. End if.
\end{enumerate}
End for.
\item Set $i=i+1$.
\end{enumerate}
\textbf{end while}. Set $\Candidates=\PC{N}$.
\paragraph{Part 2: Finding an optimal pair within Candidates}
For all pairs $g,h\in\Candidates$ and for all gates $O\in\{\land,\lor\}$ do: If $L(g)+L(h)=L(f)$ and $f=gOh$ then return $\{g,O,h\}$. End if. End for.\par
\textbf{end procedure}
\end{quote}
\subsection*{5.1 Correctness of Algorithm 2}
In this subsection we show that Algorithm 2 has the desired input/output behavior. Fix some function $f$ with $n$-inputs satisfying $L(f)\geq 2$. Let $N=2^n$.
\paragraph{Part 1: building Candidates}
First, we will prove some loop invariants that will help us show that $\Candidates$ and $\PC{i}$ contain those functions we are interested in and do not contain many more things.
The following claim shows that the $x^?$ described on Line 10 always exists and that the $?$-values of partial functions in $\PC{i}$ always have an easily computable structure.
\begin{ntclaim}[Author Claim 35]
Before and after each iteration of the while loop, it is true that if $\gamma\in\PC{i}$, then
\[\gamma(x)=?\quad\Longleftrightarrow\quad x\geq i\]
(interpreting $i$ as a binary string in $\{0,1\}^n$ in the natural way), and consequently $|\gamma^{-1}(\{0,1\})|=i$.
\end{ntclaim}
\begin{proof}
Clearly the claim is satisfied before the first iteration of the while loop when $i=0$ and $\PC{i}=\{\operatorname{AllUnknown}\}$.
Now, we must argue inductively. Suppose $1\leq i\leq N$ and $\gamma'\in\PC{i}$. Then, it follows that there is some $\gamma\in\PC{i-1}$ and some $b\in\{0,1\}$ such that $\gamma'=\gamma_b$ where $\gamma_b$ is as defined in the pseudocode. That is, $\gamma_b$ is equal to $\gamma$ except that the first $?$-value (which occurs at $x^?_{\mathrm{old}}=i-1$ by the inductive hypothesis) is replaced by a $b$. Thus, we have
\[\gamma'(x)=?\Longleftrightarrow\gamma(x)=?\land(x\ne x^?_{\mathrm{old}})\Longleftrightarrow x>x^?_{\mathrm{old}}\Longleftrightarrow x\geq i\]
where the first equivalence comes from the definition of $\gamma_b=\gamma'$ and the second equivalence comes from the fact that $x^?_{\mathrm{old}}=i-1$.
\end{proof}
Next, we show that the $\PC{i}$ never contains ``redundant'' partial functions.
\begin{ntclaim}[Author Claim 36]
Before and after each iteration of the while loop, it is true that if $\gamma'$ and $\gamma''$ are distinct elements of $\PC{i}$, then no total function agrees with both $\gamma'$ and $\gamma''$.
\end{ntclaim}
\begin{proof}
Before the first iteration of the while loop runs, $i=0$ and $\PC{0}$ only contains the single partial function $\operatorname{AllUnknown}$, so the claim clearly holds.
Now we must show that the claim holds inductively. Assume $1\leq i\leq N$. For contradiction, suppose there was some total function $q$ that agrees with distinct elements $\mu$ and $\mu'$ from $\PC{i}$. It follows that there exists some $b,b'\in\{0,1\}$ and some (possibly not distinct) $\gamma,\gamma'\in\PC{i-1}$ such that $\mu=\gamma_b$ and $\mu'=\gamma'_{b'}$ (using the notation from the pseudocode where these functions $\gamma_b$ and $\gamma'_{b'}$ are given by replacing the output of the first $?$-valued input in $\gamma$ or $\gamma'$ respectively with a $b$-value or $b'$-value respectively). It follows that $q$ must also agree with $\gamma$ and $\gamma'$.
Either $\gamma\ne\gamma'$ or not. If $\gamma\ne\gamma'$, then $q$ agrees with two distinct elements from $\PC{i-1}$ which contradicts the inductive hypothesis.
Now suppose that $\gamma=\gamma'$. Then it must be that $b\ne b'$ (otherwise, $\mu=\mu'$ and we assumed they are distinct). But then, we have then $\gamma$ and $\gamma'$ have the same first $?$-valued input $x^?$, so
\[b=\mu(x^?)=q(x^?)=\mu'(x^?)=b'\]
which contradicts that $b\ne b'$.
\end{proof}
Moreover, $\PC{i}$ only contains partial functions that can be completed to total functions that pass a certain test.
\begin{ntclaim}[Author Claim 37]
Before and after each iteration of the while loop, it is true that if $\gamma\in\PC{i}$ then there exists a function $g$ on $n$-inputs that agrees with $\gamma$ such that
\[L(\Select[f,g])\leq L(f)+n+2.\]
\end{ntclaim}
\begin{proof}
Before the first iteration of the while loop runs, $i=0$ and $\PC{0}$ only contains one partial function ($\operatorname{AllUnknown}$). The function $f$ clearly agrees with $\operatorname{AllUnknown}$, and it is easy to see that $L(\Select[f,f])=L(f)\leq L(f)+n+1$, as desired. Thus, the claim holds before the first iteration of the while loop.
Moreover, the claim clearly continues holding inductively because before any $\gamma_b$ is added to $\PC{i}$, we check to see if the function $g_{\gamma_b}$ satisfies $L(\Select[f,g_{\gamma_b}])\leq L(f)+n+2$ and $g_{\gamma_b}$ agrees with $\gamma_b$ by construction.
\end{proof}
Finally, we show that $\PC{i}$ always contains the partial functions we want.
\begin{ntclaim}[Author Claim 38]
Suppose some function $q$ is in an optimal subcomputation for $f$. Then before and after each iteration of the while loop there is a $\gamma\in\PC{i}$ such that $q$ agrees with $\gamma$. Moreover, once part 1 is finished, $q\in\Candidates$.
\end{ntclaim}
\begin{proof}
Fix some $q$ as in the statement of the claim.
Before the first iteration of the while loop runs, $i=0$ and $\PC{0}$ contains the all-$?$ partial function $\operatorname{AllUnknown}$, so $q$ agrees with $\operatorname{AllUnknown}$ and the claim holds.
Now, we must show the claim holds inductively. Assume $1\leq i\leq N$. Then by induction there exists a $\gamma\in\PC{i-1}$ such that $q$ agrees with $\gamma$. Let $b=q(i-1)$. Then $q$ agrees with $\gamma_b$ as defined in the pseudocode (replacing the first $?$-value in $\gamma$ with a $b$-value) since Claim 35 implies that
\[\gamma_b(x)=\begin{cases}b,&\text{if }x=i-1,\\\gamma(x),&\text{otherwise}.\end{cases}\]
Thus, if we could show $\gamma_b\in\PC{i}$, we would be done with showing the first part of the claim. From the pseudocode, it is clear $\gamma_b\in\PC{i}$ if
\[L(\Select[f,g_{\gamma_b}])\leq L(f)+n+2,\]
where $g_{\gamma_b}$ is as defined in the code (the function given by replacing the $?$-values in $\gamma_b$ with ones) which we now prove.
We already noted that $\gamma_b(x)=\begin{cases}b,&\text{if }x=i-1,\\\gamma(x),&\text{otherwise}.\end{cases}$ Thus, appealing to Claim 35, we know that $\gamma_b(x)=?\Longleftrightarrow x>x^?$ where $x^?\in\{0,1\}^n$ is the binary string equivalent to $i-1$ (note that $0\leq i-1\leq N-1$ so this makes sense). Hence, since $q$ agrees with $\gamma_b$, we have that $g_{\gamma_b}(x)=q(x)\lor\operatorname{GrtrThan}_{x^?}(x)$ where $\operatorname{GrtrThan}_{x^?}(x)=1$ if and only if $x>x^?$.
Thus, we have that
\begin{align*}\Select[f,g_{\gamma_b}](x,z)&=\begin{cases}f(x),&\text{if }z=0,\\g_{\gamma_b}(x),&\text{if }z=1\end{cases}\\&=\Select[f,q](x,z)\lor(z\land\operatorname{GrtrThan}_{x^?}(x)).\end{align*}
Since $\{g,O,h\}\in\Opt(f)$, we know that $L(\Select[f,q])=L(f)+1$ by Lemma 23, and Proposition 15 implies that $L(\operatorname{GrtrThan}_{x^?})\leq n$. Hence, we have that $L(\Select[f,g_{\gamma_b}])\leq L(f)+n+2$.
Finally, we show that $q\in\Candidates$ after part 1 finishes. Clearly, it suffices to show that $q\in\PC{N}$ after part 1 finishes. We have already shown that there is a $\gamma\in\PC{N}$ such that $\gamma$ agrees with $q$. However, Claim 35 implies that $\gamma$ is a total function and hence it equals $q$, so $q\in\PC{N}$.
\end{proof}
